% Options for packages loaded elsewhere
\PassOptionsToPackage{unicode}{hyperref}
\PassOptionsToPackage{hyphens}{url}
\PassOptionsToPackage{dvipsnames,svgnames,x11names}{xcolor}
\documentclass[
  10pt,
  a4paper,
]{article}
\usepackage{xcolor}
\usepackage[margin=25mm]{geometry}
\usepackage{amsmath,amssymb}
\setcounter{secnumdepth}{2}
\usepackage{iftex}
\ifPDFTeX
  \usepackage[T1]{fontenc}
  \usepackage[utf8]{inputenc}
  \usepackage{textcomp} % provide euro and other symbols
\else % if luatex or xetex
  \usepackage{unicode-math} % this also loads fontspec
  \defaultfontfeatures{Scale=MatchLowercase}
  \defaultfontfeatures[\rmfamily]{Ligatures=TeX,Scale=1}
\fi
\usepackage{lmodern}
\ifPDFTeX\else
  % xetex/luatex font selection
\fi
% Use upquote if available, for straight quotes in verbatim environments
\IfFileExists{upquote.sty}{\usepackage{upquote}}{}
\IfFileExists{microtype.sty}{% use microtype if available
  \usepackage[]{microtype}
  \UseMicrotypeSet[protrusion]{basicmath} % disable protrusion for tt fonts
}{}
\makeatletter
\@ifundefined{KOMAClassName}{% if non-KOMA class
  \IfFileExists{parskip.sty}{%
    \usepackage{parskip}
  }{% else
    \setlength{\parindent}{0pt}
    \setlength{\parskip}{6pt plus 2pt minus 1pt}}
}{% if KOMA class
  \KOMAoptions{parskip=half}}
\makeatother
\usepackage{longtable,booktabs,array}
\usepackage{caption}
\captionsetup[table]{skip=6pt}
\newcounter{none} % for unnumbered tables
\usepackage{calc} % for calculating minipage widths
% Correct order of tables after \paragraph or \subparagraph
\usepackage{etoolbox}
\makeatletter
\patchcmd\longtable{\par}{\if@noskipsec\mbox{}\fi\par}{}{}
\makeatother
% Allow footnotes in longtable head/foot
\IfFileExists{footnotehyper.sty}{\usepackage{footnotehyper}}{\usepackage{footnote}}
\makesavenoteenv{longtable}
\setlength{\emergencystretch}{3em} % prevent overfull lines
\providecommand{\tightlist}{%
  \setlength{\itemsep}{0pt}\setlength{\parskip}{0pt}}
\usepackage{amsmath,amssymb,booktabs,longtable,array,microtype}
\usepackage{fancyhdr}
\setlength{\headheight}{14pt}
\pagestyle{fancy}
\fancyhf{}
\fancyhead[L]{\small C600 Studio -- Puzzle Theory}
\fancyhead[R]{\small\thepage}
\renewcommand{\headrulewidth}{0.3pt}
\setlength{\emergencystretch}{3em}
\linespread{1.06}
\renewcommand{\arraystretch}{1.16}
\setcounter{tocdepth}{2}
\AtBeginDocument{\hypersetup{pdfauthor={C600 Studio Research Supplement},
  pdftitle={Puzzle Theory of the Full 600-Cell},
  pdfsubject={Geometry, topology, group actions, and sound optimization},
  pdfkeywords={600-cell, MPUlt, H4, grip theory, Rethlas, finite permutation groups}}}
\newcommand{\Ico}{I^{*}}
\newcommand{\supp}{\operatorname{supp}}
\newcommand{\Stab}{\operatorname{Stab}}
\newcommand{\conv}{\operatorname{conv}}
\newcommand{\Isom}{\operatorname{Isom}}
\newcommand{\sgn}{\operatorname{sgn}}
\usepackage{bookmark}
\IfFileExists{xurl.sty}{\usepackage{xurl}}{} % add URL line breaks if available
\urlstyle{same}
% fallback for those not using the hyperref driver hyperxmp:
\makeatletter
\@ifundefined{xmpquote}{\newcommand{\xmpquote}[1]{#1}}{}
\makeatother
\hypersetup{
  pdftitle={Puzzle Theory of the Full 600-Cell},
  pdfauthor={C600 Studio Research Supplement},
  colorlinks=true,
  linkcolor={blue},
  filecolor={Maroon},
  citecolor={Blue},
  urlcolor={blue},
  pdfcreator={LaTeX via pandoc}}

\title{Puzzle Theory of the Full 600-Cell}
\usepackage{etoolbox}
\makeatletter
\providecommand{\subtitle}[1]{% add subtitle to \maketitle
  \apptocmd{\@title}{\par {\large #1 \par}}{}{}
}
\makeatother
\subtitle{Geometry, Topology, Group Actions, and Sound Optimization}
\author{C600 Studio Research Supplement}
\date{7 September 2026}

\begin{document}
\maketitle

{
\setcounter{tocdepth}{2}
\tableofcontents
}\clearpage
\section{Purpose and evidence}\label{purpose-and-evidence}

The full-cut 600-cell makes a useful meeting point between
regular-polytope geometry, finite group actions, constructive
algorithms, and graphics engineering. Its exterior has only 600
tetrahedral cells, but the retained subdivision has 259,800 sticker
slots on 177,120 surface pieces. Consequently, understanding the base
polytope is necessary but insufficient for understanding the puzzle. A
drawing of the 600-cell, a labelled puzzle configuration, and a camera
pose are three different mathematical objects.

This supplement develops that distinction systematically. The geometry
determines incidence, local symmetry and the topology of the supporting
space. The cuts determine which regions move together. Legal
permutations determine reachability and invariants. Ordered local frames
make orientation explicit, while complete permutation witnesses connect
mathematical controllers to recoverable software operations. The final
engineering section explains which optimizations preserve those
contracts and which need further measurement or proof.

The principal software contribution remains a modification of
\textbf{Magic Puzzle Ultimate by Andrey Astrelin}, who receives primary
credit for the original simulator and renderer. Nan Ma's solving
methodologies and ivan216's projects are background to the broader
development history, as documented in the companion technical report.
They are not attributed the new derivations or the Rethlas review in
this supplement.

\subsection{Three levels of evidence}\label{three-levels-of-evidence}

{\def\LTcaptype{none} % do not increment counter
\begin{longtable}[]{@{}
  >{\raggedright\arraybackslash}p{(\linewidth - 4\tabcolsep) * \real{0.3333}}
  >{\raggedright\arraybackslash}p{(\linewidth - 4\tabcolsep) * \real{0.3333}}
  >{\raggedright\arraybackslash}p{(\linewidth - 4\tabcolsep) * \real{0.3333}}@{}}
\toprule\noalign{}
\begin{minipage}[b]{\linewidth}\raggedright
Level
\end{minipage} & \begin{minipage}[b]{\linewidth}\raggedright
What is established
\end{minipage} & \begin{minipage}[b]{\linewidth}\raggedright
Boundary
\end{minipage} \\
\midrule\noalign{}
\endhead
\bottomrule\noalign{}
\endlastfoot
Exact base geometry & Coordinate arithmetic, incidence, quaternion
closure, explicit symmetries, mod-2 boundary ranks and graph distances &
Does not reconstruct the full cut subdivision \\
Retained puzzle certificates & Orbit census, legal generator maps, local
frame groups, parity constraints and selected controller certificates &
Conditional inputs to the theorems; not an exact legal-group order
computation \\
Engineering analysis & Consequences for state updates, search,
rendering, picking and qualification experiments & Proposals are not
claimed as implemented speedups or new GPU benchmarks \\
\end{longtable}
}

The geometry checker uses integer pairs in \(\mathbb Z[\phi]\), with
\(\phi^2=\phi+1\). It does not infer incidence from floating-point
proximity. A separate topological argument identifies the boundary with
\(S^3\); matching homology ranks alone would not prove that
identification. The mathematical blueprint that follows was generated
and independently reviewed using Rethlas with a documented Windows I/O
adaptation. Its review is an informal AI proof check, not a Lean proof
or peer review. The retained artifacts record the actual verdict and
scope.

\subsection{How the HSC wiki is used}\label{how-the-hsc-wiki-is-used}

The Hypercubing wiki associated with the Hyperspeedcube community
supplies useful definitions and examples. Its \textbf{Formal
Introduction to Grip Theory} distinguishes solved grips, current active
grips and attitude; this helps connect geometric cap membership to piece
motion. Its \textbf{Piece Invariants} explains how parity and abelian
orientation information constrain reachable states. Its \textbf{God's
Number} emphasizes a specified move metric when deriving counting
bounds, and \textbf{Cut Depths} provides normalized cut-distance
conventions.

Those pages do not supply this full-cut 600-cell's 35-orbit census. In
particular, the wiki's \(A_4\) monoflip example belongs to a hypercube
setting; the full 600-cell's twenty-sticker moving orbit has the
retained local orientation group \(A_5\). The relevant wiki pages are
linked in the reading guide. Their definitions guide the exposition; the
specialized arguments and finite checks are provided here.

The term \emph{complete description} below means a connected account of
the specified model and the limits of present knowledge. It does not
assert a presentation or exact order for its entire legal twist group, a
classification of all cut-depth regimes, or an optimal solver.

\subsection{Objects that must remain
distinct}\label{objects-that-must-remain-distinct}

{\def\LTcaptype{none} % do not increment counter
\begin{longtable}[]{@{}
  >{\raggedright\arraybackslash}p{(\linewidth - 4\tabcolsep) * \real{0.3333}}
  >{\raggedright\arraybackslash}p{(\linewidth - 4\tabcolsep) * \real{0.3333}}
  >{\raggedright\arraybackslash}p{(\linewidth - 4\tabcolsep) * \real{0.3333}}@{}}
\toprule\noalign{}
\begin{minipage}[b]{\linewidth}\raggedright
Object
\end{minipage} & \begin{minipage}[b]{\linewidth}\raggedright
Size or type
\end{minipage} & \begin{minipage}[b]{\linewidth}\raggedright
Role
\end{minipage} \\
\midrule\noalign{}
\endhead
\bottomrule\noalign{}
\endlastfoot
\(P\) & Convex four-dimensional polytope & Geometric body supporting the
construction \\
\(\partial P\) & Topological three-sphere & Exterior space subdivided
into tetrahedra and sticker regions \\
\(I^*\) & 120 unit quaternions & Coordinate group underlying the regular
vertices \\
\(K^+\) & 7,200 elements & Orientation-preserving rigid symmetries of
the whole polytope \\
\(K=W(H_4)\) & 14,400 elements & Full rigid symmetry group, including
reflections \\
\(SO(4)\) & Continuous group, dimension six & Intended 4D camera
rotations \\
Cap rotation group & \(A_4\), order 12 & Rotations permitted within one
tetrahedral cap \\
\(H_o\) & Retained local frame group & Orientations of one piece in
moving orbit \(o\) \\
\(G\) & Finite; exact order not established here & Group generated by
the legal twists on labelled sticker slots \\
\end{longtable}
}

Topological dimension is equally important. A tetrahedral cell of the
600-cell is three-dimensional, and a sticker region in that cell has
three-dimensional relative interior. Its displayed mesh uses triangles
to draw two-dimensional boundary faces. The application's rank is a
cap-membership statistic and may exceed four; it is not an ambient
dimension. A legal permutation of separated sticker regions is also not
automatically a continuous self-map of the originally glued boundary. In
particular, the trivial fundamental group of \(S^3\) says nothing about
the number of legal scrambles or the noncommutativity of twists.

\newpage

\section{Mathematical development}\label{mathematical-development}

The following blueprint separates unconditional statements about the
regular base polytope from implications that explicitly assume the
retained puzzle certificates. Products of moves use the stated
chronological convention. A geometric symmetry is not silently counted
as a free puzzle move.

\subsection{Definitions and mathematical
sources}\label{definitions-and-mathematical-sources}

Put \begin{equation}
\phi=\frac{1+\sqrt5}{2},\qquad \tau=\phi^{-1}=\phi-1.
\end{equation} Thus \begin{equation}
\phi^2=\phi+1,\qquad \tau^2=2-\phi.
\end{equation}

Identify \(\mathbb R^4\) with the quaternion algebra \(\mathbb H\) using
the ordered orthonormal basis \(1,i,j,k\). Quaternion conjugation is
denoted by \(x\mapsto\bar x\), and \begin{equation}
\operatorname{Sp}(1)=\{x\in\mathbb H:|x|=1\}.
\end{equation} The binary icosahedral group \(I^*\) is the inverse image
of the rotation group of a regular icosahedron under \begin{equation}
\operatorname{Sp}(1)\longrightarrow SO(3),\qquad
p\longmapsto\bigl(u\mapsto pup^{-1}\bigr),
\end{equation} where \(\mathbb R^3\) is identified with the imaginary
quaternions.

For a finite set \(X\), write \(S_X\) and \(A_X\) for its symmetric and
alternating groups. Write \(C_m\) for the cyclic group of order \(m\),
and \begin{equation}
D_5=\langle R,F\mid R^5=F^2=1,\ FRF=R^{-1}\rangle
\end{equation} for the dihedral group of order \(10\).

A flag of a four-dimensional polytope is an incident chain consisting of
a vertex, edge, two-face, and facet. A polytope is regular when its
geometric symmetry group is transitive on flags.

Two external geometric results will be used, with the following complete
statements and conventions.

\textbf{Quaternion coordinate result.} In a suitable orthonormal
quaternion coordinate system, the inverse image of the icosahedron's
rotation group consists precisely of \begin{equation}
(\pm1,0,0,0),\qquad
\frac12(\pm1,\pm1,\pm1,\pm1),\qquad
\frac12(\pm\phi,\pm1,\pm\tau,0),
\end{equation} together with their even coordinate permutations, with
independent signs.

Source identifiers:
\texttt{paper\_id\ =\ Baez,\ From\ the\ Icosahedron\ to\ E8};
\texttt{arXiv\_id\ =\ 1712.06436v2};
\texttt{theorem\_id\ =\ unnumbered\ coordinate\ assertion\ in\ “The\ Icosians”}.
Here the finite group is the binary lift, not the infinite ring of
integral combinations called the icosians.
\href{https://arxiv.org/html/1712.06436v2}{Baez, ``The Icosians''}.

\textbf{Binary-group polytope result.} The binary icosahedral subgroup
of unit quaternions is one orbit of the reflection group of type
\(H_4\), and its elements are the vertices of a regular 600-cell.

Source identifiers:
\texttt{paper\_id\ =\ Choi-\/-Lee,\ Binary\ Icosahedral\ Group\ and\ 600-Cell,\ Symmetry\ 10\ (2018),\ 326};
\texttt{theorem\_id\ =\ Theorem\ 2};
\texttt{arXiv\_id\ =\ not\ supplied}. The paper's binary-group notation
indicates a double covering, not coordinate dilation. Its proof
establishes reflection invariance in Lemma 2, then identifies the orbit
of \(1\) using its \(H_3\) stabilizer and the \(120\)-element orbit
count. Its quaternion norm is the ordinary Euclidean norm used here.
\href{https://pdfs.semanticscholar.org/0fef/c9c8c31093da12262c8f6740da46c0dfb46e.pdf}{Choi--Lee,
Theorem 2 and its proof}.

These results concern the uncut polytope. They provide no sticker census
or puzzle reachability theorem.

\subsection{Lemma 1: Coordinate
identification}\label{lemma-1-coordinate-identification}

\subsubsection{Statement}\label{statement}

The set \(V\) in the problem has \(120\) distinct elements, every
element has norm \(2\), and \(V/2\) is a binary icosahedral subgroup of
\(\operatorname{Sp}(1)\). Its convex hull is a regular 600-cell.

\subsubsection{Proof}\label{proof}

The three coordinate families are disjoint because their patterns of
zero coordinates and absolute values differ. Their cardinalities are
\begin{equation}
8,\qquad 16,\qquad 12\cdot8=96.
\end{equation} In the third family all four absolute values are
distinct, so the twelve even permutations give twelve distinct
arrangements for each sign choice.

The squared norms are respectively \begin{equation}
4,\qquad 4,\qquad 1+\phi^2+\tau^2=4.
\end{equation}

The permutation taking \begin{equation}
(\phi,1,\tau,0)\quad\text{to}\quad(0,1,\phi,\tau)
\end{equation} has index sequence \((4,2,1,3)\) and is even. Thus the
displayed coordinate set agrees exactly with twice the quaternion
coordinate set stated above, including its permutation convention.

The coordinate result identifies \(V/2\) with \(I^*\). This
identification uses the inverse image of an icosahedral rotation group,
hence supplies multiplication closure and inverses, not merely a set of
\(120\) points on a sphere. One direct route to its coordinates is to
lift the identity and the rotations about edge, face, and vertex axes
using \begin{equation}
\cos(\theta/2)+u\sin(\theta/2)
\end{equation} and its negative, where \(u\) is the unit imaginary axis.
The rotations comprise the identity, fifteen half-turns, twenty
third-turns, and twenty-four fifth-turn rotations.

The binary-group polytope result now applies in the same unit-quaternion
metric. Multiplication by \(2\) preserves regularity and the face
lattice. Therefore \(P=\operatorname{conv}(V)\) is a regular 600-cell.
\(\square\)

\subsection{Lemma 2: Radial topology of the
boundary}\label{lemma-2-radial-topology-of-the-boundary}

\subsubsection{Statement}\label{statement-1}

The boundary \(\partial P\) is homeomorphic to \(S^3\). Consequently,
\begin{equation}
H_j(\partial P;\mathbb Z)\cong
\begin{cases}
\mathbb Z,&j=0,3,\\
0,&j\ne0,3,
\end{cases}
\qquad
\pi_1(\partial P)=1.
\end{equation}

\subsubsection{Proof}\label{proof-1}

The polytope contains \begin{equation}
\operatorname{conv}\{\pm2e_1,\ldots,\pm2e_4\},
\end{equation} which contains the Euclidean unit ball. Indeed,
\(\|x\|_2\le1\) implies \(\|x\|_1\le2\). Thus \(0\) lies in the interior
of \(P\). The polytope is compact and contained in the ball of radius
\(2\).

For every \(u\in S^3\), convexity and compactness give \begin{equation}
P\cap\{tu:t\ge0\}=\{tu:0\le t\le r(u)\}
\end{equation} for a unique \(r(u)>0\). Every point \(tu\) with
\(0\le t<r(u)\) is interior: it is a convex combination, with positive
weight on the first term, of an interior point \(0\) and the point
\(r(u)u\in P\). Hence \(r(u)u\) is the unique boundary point on this
ray.

The map \begin{equation}
\partial P\longrightarrow S^3,\qquad x\longmapsto\frac{x}{\|x\|}
\end{equation} is therefore a continuous bijection. Its domain is
compact and its codomain Hausdorff, so it is a homeomorphism.

The sphere \(S^3\) has a CW decomposition with one zero-cell and one
three-cell, giving the stated integral homology. For simple
connectedness, cover \(S^3\) by the complements of two distinct
antipodal points. Each complement is homeomorphic to \(\mathbb R^3\),
and their intersection is path connected. The van Kampen theorem gives
trivial fundamental group.

This proves the topology directly. Neither mod-\(2\) Betti numbers nor
equality of integral homology groups alone would establish the
homeomorphism or fundamental group. \(\square\)

\subsection{Lemma 3: Faces, links, metric data, and
duality}\label{lemma-3-faces-links-metric-data-and-duality}

\subsubsection{Statement}\label{statement-2}

The polytope \(P\) has circumradius \(2\), edge length \(2/\phi\), and
\begin{equation}
(f_0,f_1,f_2,f_3)=(120,720,1200,600).
\end{equation} Its facets are regular tetrahedra, five facets meet at
each edge, twenty meet at each vertex, and its vertex links are
icosahedral. Its Schläfli symbol is \(\{3,3,5\}\), and its polar dual is
a regular 120-cell of type \(\{5,3,3\}\).

\subsubsection{Proof}\label{proof-2}

Every point of \(V\) is an exposed vertex: the functional
\(x\mapsto v\cdot x\) has its unique maximum at \(v\) among points of
\(V\), since all have norm \(2\). Thus \(P\) has exactly \(120\)
vertices and circumradius \(2\).

Let \(v=(2,0,0,0)\). The largest first coordinate of another vertex is
\(\phi\), attained at twelve vertices. Consequently the largest inner
product with \(v\) among distinct vertices is \(2\phi\), and the minimum
squared distance is \begin{equation}
8-4\phi=4\tau^2.
\end{equation} The group \(I^*\) acts transitively on \(V\) by left
multiplication, so this is also the global minimum distance.

A closest pair of equal-norm vertices is an edge. To see this, let their
common squared norm be \(R^2\) and their inner product be \(m<R^2\). For
any other vertex \(z\), minimality of their distance gives
\begin{equation}
(v+w)\cdot z\le2m<R^2+m=(v+w)\cdot v=(v+w)\cdot w.
\end{equation} Thus the functional \(v+w\) exposes precisely their
segment.

Regularity makes all edges congruent. Hence the twelve closest vertices
are exactly the neighbors of \(v\), and the edge length is
\(2\tau=2/\phi\).

Their projections onto \(v^\perp\) have coordinates \begin{equation}
(0,\pm1,\pm\tau)
\end{equation} and cyclic permutations. These are a scaled, orthogonally
positioned regular icosahedron. All twelve neighboring vertices have
first coordinate \(\phi\), so a sufficiently close hyperplane section at
\(v\) is similar to this projected set. It follows that the vertex
figure, and hence the vertex link, is icosahedral.

An icosahedron has twelve vertices, thirty edges, and twenty triangular
faces. These numbers can also be obtained directly from its coordinates:
each vertex has five nearest neighbors arranged in a pentagon; the
resulting triangular boundary has \begin{equation}
E=12\cdot5/2=30,\qquad F=2E/3=20.
\end{equation} Its triangular faces correspond to facets through \(v\),
giving twenty tetrahedral facets through each vertex. Tetrahedrality
also follows from the regular 600-cell identification in Lemma 1.

Double counting yields \begin{equation}
f_1=\frac{120\cdot12}{2}=720,
\qquad
f_3=\frac{120\cdot20}{4}=600.
\end{equation} Each tetrahedron has four triangular faces, and each
ridge of a convex four-polytope belongs to two facets. Therefore
\begin{equation}
f_2=\frac{600\cdot4}{2}=1200.
\end{equation} The five triangles incident to a vertex of an icosahedral
link correspond to the five facets incident to the associated edge of
\(P\). Thus the Schläfli symbol is \(\{3,3,5\}\).

The Euler characteristic is \begin{equation}
120-720+1200-600=0,
\end{equation} in agreement with Lemma 2.

Because \(0\) is interior, the polar \begin{equation}
P^\circ=\{y:y\cdot x\le1\text{ for every }x\in P\}
\end{equation} is a bounded four-polytope. Polarity reverses proper face
incidences and sends a \(j\)-face to a \((3-j)\)-face. Orthogonal
symmetries commute with polarity, so the dual is regular. Its facets are
dual to the icosahedral vertex figures and hence are dodecahedra. It has
\begin{equation}
(f_0,f_1,f_2,f_3)=(600,1200,720,120)
\end{equation} and reversed symbol \(\{5,3,3\}\): it is the regular
120-cell.

This duality is not an identification of the two polytopes or their face
counts. \(\square\)

\subsection{Lemma 4: The finite geometric symmetry
group}\label{lemma-4-the-finite-geometric-symmetry-group}

\subsubsection{Statement}\label{statement-3}

For \(p,q\in I^*\), define \begin{equation}
L_{p,q}(x)=pxq^{-1}.
\end{equation} Then \begin{equation}
K^+\cong (I^*\times I^*)/\langle(-1,-1)\rangle,
\qquad |K^+|=7200,
\end{equation} and quaternion conjugation extends this to \(K\) of order
\(14400\).

Moreover, \begin{equation}
K\cong W(H_4)
\end{equation} with presentation \begin{equation}
\left\langle s_0,s_1,s_2,s_3\ \middle|\
\begin{array}{l}
s_i^2=1,\\
(s_0s_1)^3=(s_1s_2)^3=(s_2s_3)^5=1,\\
(s_is_j)^2=1\quad\text{if }|i-j|>1
\end{array}
\right\rangle.
\end{equation}

\subsubsection{Proof}\label{proof-3}

Quaternion norm multiplicativity makes \(L_{p,q}\) orthogonal. Its
determinant is positive: the same construction on the connected space
\(\operatorname{Sp}(1)\times\operatorname{Sp}(1)\) depends continuously
on \((p,q)\) and includes the identity. Multiplication closure of
\(I^*\) shows that \(L_{p,q}\) preserves \(V/2\), and therefore \(P\).

If \(L_{p,q}\) is the identity, evaluating at \(1\) gives \(p=q\). This
quaternion must then commute with every quaternion. Commuting with \(i\)
and \(j\) forces it to be real, so it is \(1\) or \(-1\). The kernel is
precisely \begin{equation}
\{(1,1),(-1,-1)\}.
\end{equation} Thus these maps give \(120^2/2=7200\) distinct
positive-determinant symmetries.

Quaternion conjugation \begin{equation}
J(x)=\bar x
\end{equation} preserves \(V\), since signs are independent. Its real
matrix is \begin{equation}
\operatorname{diag}(1,-1,-1,-1),
\end{equation} so \(\det J=-1\). The maps \(L_{p,q}\) and \(L_{p,q}J\)
therefore give \(14400\) distinct symmetries.

An upper bound is still necessary. Fix a vertex \(v\). Its stabilizer
acts faithfully on the icosahedral vertex figure: the vertex direction
and the three-dimensional span of that figure together span
\(\mathbb R^4\). An orthogonal symmetry of an icosahedron is determined
by the image of one ordered triangular face. There are at most
\begin{equation}
20\cdot6=120
\end{equation} such images. Orbit--stabilizer consequently gives
\begin{equation}
|K|\le120\cdot120=14400.
\end{equation} The exhibited maps attain this upper bound, proving both
group orders and the description of \(K^+\).

Also, \begin{equation}
J L_{p,q}J=L_{q,p},
\end{equation} so conjugation exchanges the two quaternion factors.

For the Coxeter presentation, radially project the barycentric
subdivision of \(\partial P\) onto \(S^3\). A flag gives a spherical
tetrahedral chamber. A symmetry fixing a flag fixes the barycenters of
its four faces; these four vectors are linearly independent, so the
symmetry is the identity. Regularity therefore gives a simply transitive
action on chambers.

For each of the four chamber walls, the symmetry taking the chosen
chamber to its adjacent chamber fixes the other three barycenter rays.
It is the reflection in their span. Denote these reflections by
\(s_0,\ldots,s_3\), indexed by face ranks.

Chamber adjacency is connected, so these reflections generate \(K\).
Consecutive rank changes give rank-two sections of sizes \(3,3,5\),
producing the displayed relations. Nonconsecutive changes commute and
give square galleries.

These relations are complete. A word acting trivially determines a
closed gallery in the chamber complex. The dual cell decomposition is a
cell decomposition of \(S^3\), which is simply connected by Lemma 2. Its
two-skeleton therefore has trivial fundamental group. Closed galleries
are generated by backtracking and boundaries of dual two-cells. Those
boundaries are precisely the involution, square, and rank-two polygon
relations above. Hence the displayed presentation presents \(K\), rather
than merely mapping onto it.

The Coxeter group with this chain presentation is \(W(H_4)\).
\(\square\)

\subsection{The finite checker and its
scope}\label{the-finite-checker-and-its-scope}

The supplied checker implements a finite audit in the exact ring
\begin{equation}
\mathbb Z[\phi]=\{a+b\phi:a,b\in\mathbb Z\}.
\end{equation} It represents an element by \((a,b)\) and uses
\begin{equation}
(a+b\phi)(c+d\phi)
=(ac+bd)+(ad+bc+bd)\phi.
\end{equation} Its sign test compares integers after writing
\begin{equation}
2(a+b\phi)=(2a+b)+b\sqrt5,
\end{equation} with sign cases handled before squaring. Thus its
incidence tests do not depend on floating-point tolerances.

Its algorithm:

\begin{enumerate}
\def\labelenumi{\arabic{enumi}.}
\tightlist
\item
  Enumerates the three vertex families and verifies norms.
\item
  Joins pairs with inner product \(2\phi\).
\item
  Enumerates graph triangles and four-cliques.
\item
  For each four-clique, forms the sum \(N\) of its vertices and checks a
  strict supporting inequality at every other vertex.
\item
  Counts incidences and constructs vertex links.
\item
  Computes boundary matrices and ranks over \(\mathbb F_2\).
\item
  Tests quaternion multiplication closure and counts distinct linear
  maps by their images on a basis.
\end{enumerate}

For a tetrahedron with vertices \(v_1,\ldots,v_4\), the Gram matrix has
diagonal entries \(4\) and off-diagonal entries \(2\phi\). It is
positive definite. Moreover, \begin{equation}
N=\sum_i v_i,\qquad
N\cdot v_i=4+6\phi,\qquad
\|N\|^2=16+24\phi.
\end{equation} Thus the supporting tests certify actual tetrahedral
facets, not merely graph cliques.

There is also a completeness argument beyond finding some supporting
facets. If every triangular ridge of the listed tetrahedra has both its
incident facets in the list, the list is closed under facet adjacency.
The facet adjacency graph of a convex polytope is connected, so a
nonempty such list contains all facets.

Nevertheless:

\begin{itemize}
\tightlist
\item
  A link's \(f\)-vector alone does not identify it as an icosahedron;
  Lemma 3 supplies its geometry.
\item
  Ranks over \(\mathbb F_2\) do not prove integral homology or simple
  connectedness; Lemma 2 supplies those conclusions.
\item
  Exhibiting \(14400\) maps does not exclude further symmetries; Lemma 4
  supplies the upper bound.
\item
  No part of this checker reconstructs the cut subdivision, verifies all
  legal sticker maps, or determines the full legal group.
\end{itemize}

The supplied JSON records a previous reported run. The unconditional
proofs above do not take its \texttt{passed} field as a substitute for
these arguments.

\subsection{Lemma 5: Cell-pole
normalization}\label{lemma-5-cell-pole-normalization}

\subsubsection{Statement}\label{statement-4}

The radius-two model has inradius \begin{equation}
r_{\mathrm{in}}=\frac{\phi^2}{\sqrt2}.
\end{equation} A similarity with scale factor \(4/\phi\) and a suitable
rotation places a chosen facet pole at \begin{equation}
n_0=(\phi^3,1,1,1),
\end{equation} with every facet plane written as \begin{equation}
n_c\cdot x=\|n_c\|^2.
\end{equation}

\subsubsection{Proof}\label{proof-4}

For a tetrahedral facet, use \(N\) from the preceding calculation. Since
\begin{equation}
\|N\|^2=4(4+6\phi),
\end{equation} the perpendicular foot of the origin on its plane is
\begin{equation}
n=\frac{N}{4}.
\end{equation} Its squared norm is \begin{equation}
\|n\|^2=\frac{4+6\phi}{4}
=\frac{\phi^4}{2}.
\end{equation} Regularity gives the same inradius at every facet.

On the other hand, \begin{equation}
\|n_0\|^2=\phi^6+3=8\phi^2.
\end{equation} Therefore \begin{equation}
\frac{\|n_0\|}{r_{\mathrm{in}}}=\frac4\phi.
\end{equation} Choose \(Q\in SO(4)\) taking the direction of one
original pole to that of \(n_0\), and put \begin{equation}
\widetilde P=\frac4\phi QP.
\end{equation} Every original pole \(n\) becomes
\(\widetilde n=(4/\phi)Qn\), and its plane becomes \begin{equation}
\widetilde n\cdot x=\|\widetilde n\|^2.
\end{equation} The selected pole is exactly \(n_0\).

All subsequent puzzle geometry uses \(\widetilde P\), identified with
\(P\) by this fixed similarity. Its geometric symmetry group is the
conjugate of \(K\) under the similarity. The radius-two and cell-pole
coordinate scales are not numerically identified. \(\square\)

\subsection{Retained-model definitions and
hypotheses}\label{retained-model-definitions-and-hypotheses}

Let \(\mathcal C\) be the set of \(600\) facets, and let \(n_c\) denote
their poles in the normalization of Lemma 5. Put \begin{equation}
\alpha=\frac{121}{125},
\qquad
U_c=\{x\in\widetilde P:n_c\cdot x>\alpha\|n_c\|^2\}.
\end{equation} Central symmetry pairs each \(n_c\) with \(-n_c\). Thus
there are \(600\) oriented caps or, equivalently, \(300\) antipodal
axes. For one axis, the two outer layers satisfy \begin{equation}
n_c\cdot x>\alpha\|n_c\|^2
\quad\text{and}\quad
n_c\cdot x<-\alpha\|n_c\|^2.
\end{equation} The middle layer lies between them and is fixed during
that axis's outer-layer turns. It is not asserted to be fixed under
turns about every other axis.

A full-dimensional cut chamber is specified by choosing a strict side of
every cut hyperplane and intersecting these choices with
\(\operatorname{int}\widetilde P\). A nonempty such intersection is
convex and hence connected. A surface piece is a retained chamber whose
closure meets the boundary in one or more relatively three-dimensional
patches. Cut-boundary strata are attached consistently to these pieces;
they are not additional sticker slots.

For a piece position \(A\), define its complete cap signature
\begin{equation}
M(A)=\{c:\text{the chamber interior of }A\text{ lies in }U_c\}.
\end{equation} Define its hosting set by \begin{equation}
\operatorname{Host}(A)
=
\{h:\overline A\cap h\text{ contains a retained relatively three-dimensional patch}\}.
\end{equation} A sticker slot is such a patch. In this full-cut
description it is addressed by \begin{equation}
(M(A),h),\qquad h\in\operatorname{Host}(A).
\end{equation}

The following hypotheses specify exactly the retained evidence used
below.

\textbf{R --- Retained subdivision and legal action.} The retained model
is the complete surface subdivision just described, with no hidden
identification of distinct pieces having the same recorded address.
Legal cap rotations carry whole pieces and their sticker patches to
whole pieces and patches. It has \(259800\) sticker slots, \(177120\)
surface pieces, \(600\) one-sticker centers fixed pointwise by every
legal move, and the following \(35\) moving orbits. The frame groups are
faithful groups of permutations of the stickers of a piece after
coherent transport.

{\def\LTcaptype{none} % do not increment counter
\begin{longtable}[]{@{}lrrrl@{}}
\toprule\noalign{}
Orbit & Rank & Stickers per piece & Pieces \(N_o\) & Frame group
\(H_o\) \\
\midrule\noalign{}
\endhead
\bottomrule\noalign{}
\endlastfoot
O00 & 1 & 2 & 1200 & \(C_2\) \\
O01 & 2 & 1 & 3600 & \(1\) \\
O02 & 3 & 1 & 2400 & \(1\) \\
O03 & 3 & 2 & 3600 & \(C_2\) \\
O04 & 4 & 1 & 7200 & \(1\) \\
O05 & 4 & 1 & 7200 & \(1\) \\
O06 & 4 & 5 & 720 & \(D_5\) \\
O07 & 5 & 1 & 7200 & \(1\) \\
O08 & 5 & 2 & 3600 & \(C_2\) \\
O09 & 6 & 2 & 7200 & \(1\) \\
O10 & 7 & 1 & 7200 & \(1\) \\
O11 & 7 & 2 & 3600 & \(C_2\) \\
O12 & 8 & 1 & 7200 & \(1\) \\
O13 & 8 & 1 & 7200 & \(1\) \\
O14 & 8 & 2 & 7200 & \(1\) \\
O15 & 9 & 1 & 2400 & \(1\) \\
O16 & 9 & 2 & 7200 & \(1\) \\
O17 & 9 & 5 & 1440 & \(C_5\) \\
O18 & 10 & 1 & 7200 & \(1\) \\
O19 & 10 & 1 & 7200 & \(1\) \\
O20 & 10 & 2 & 7200 & \(1\) \\
O21 & 11 & 1 & 7200 & \(1\) \\
O22 & 11 & 2 & 3600 & \(C_2\) \\
O23 & 12 & 2 & 7200 & \(1\) \\
O24 & 13 & 1 & 7200 & \(1\) \\
O25 & 13 & 2 & 3600 & \(C_2\) \\
O26 & 14 & 1 & 7200 & \(1\) \\
O27 & 14 & 1 & 7200 & \(1\) \\
O28 & 14 & 5 & 1440 & \(C_5\) \\
O29 & 15 & 1 & 2400 & \(1\) \\
O30 & 15 & 2 & 7200 & \(1\) \\
O31 & 16 & 1 & 7200 & \(1\) \\
O32 & 17 & 2 & 3600 & \(C_2\) \\
O33 & 18 & 1 & 2400 & \(1\) \\
O34 & 19 & 20 & 120 & \(A_5\) \\
\end{longtable}
}

These data satisfy \begin{equation}
\sum_oN_o=176520,\qquad
600+\sum_ok_oN_o=259800.
\end{equation} The \(600\) fixed centers are \(600\) singleton position
orbits, not another moving orbit.

\textbf{F --- Full local cap action.} Every cap admits its full
tetrahedral rotation group \(A_4\), acting on the same complete
subdivision. The chosen \(H_c\) is a nonidentity half-turn and \(T_c\)
is a third-turn in that action.

\textbf{I --- Retained generator-invariant certificate.} In the
transported coordinates constructed below, each chosen primitive has
even positional permutation separately in every moving orbit and zero
total orientation image in each orbit's abelianization. This is a
hypothesis about every generator, not an inference from sample
scrambles.

\textbf{C --- Nine pure-controller certificates.} For \begin{equation}
\mathcal I=\{1,4,5,7,18,19,27,31,33\},
\end{equation} each orbit has a legal word whose complete sticker
permutation is a single three-cycle in that orbit and fixes every other
sticker. A verified legal relocation produces a cycle
\((A_o\,B_o\,C_o)\). For every other position \(X\) in that orbit, a
verified legal setup fixes the two buffer sticker slots \(A_o,B_o\)
pointwise and sends \(C_o\) to \(X\). Coverage means valid transitions
and coverage of every nonbuffer destination, not merely a reported
search-node count.

The solved coloring has \(600\) distinct home-cell symbols, with six
copies of each in O01, four in O33, and twelve in each of the other
seven selected orbits.

\textbf{E --- Equivariant geometric realization.} For the extended
symmetry action, the entire retained subdivision and its piece grouping
are preserved by the geometric action of \(K\). Each center slot is
geometrically distinguished at its corresponding pole \(n_c\), and
geometric permutations act on the same labelled sticker set as the legal
moves. In particular, this is a faithful geometric realization, not an
action on a collapsed display quotient.

Hypotheses R, F, I, C, and E have different roles. In particular, C does
not certify I, and neither certifies a reconstruction of all cuts. The
supplied arithmetic and controller summaries are retained evidence for
these hypotheses; their reported pass fields do not replace their full
hypotheses.

No independent derivation of the \(433\) retained regions per
tetrahedron is claimed.

\subsection{Lemma 6: Cap actions, signatures, and piece
equivariance}\label{lemma-6-cap-actions-signatures-and-piece-equivariance}

\subsubsection{Statement}\label{statement-5}

Under R and F, \(H_c,T_c\) generate the full \(A_4\) action at each cap.
The sticker-to-piece map \begin{equation}
\rho:\mathcal S\longrightarrow\mathcal P
\end{equation} is equivariant under every legal move and hence every
legal word: \begin{equation}
\rho(g(s))=\bar g(\rho(s)).
\end{equation} Complete cap signatures identify the cut-chamber
positions in this model, and \begin{equation}
r(A)=|M(A)|-1,\qquad
k(A)=|\operatorname{Host}(A)|.
\end{equation} The quantity \(r(A)\) is neither a topological dimension
nor a codimension.

\subsubsection{Proof}\label{proof-5}

The three nonidentity half-turns of a tetrahedron form the nonidentity
elements of a normal Klein four subgroup of \(A_4\). Conjugation by any
third-turn cycles these three elements. Thus \(H_c\) and its conjugates
by \(T_c\) generate the Klein four group, and adjoining \(T_c\) gives
all twelve elements of \(A_4\).

All stickers of a physical piece travel together under a legal rigid cap
rotation by R. Therefore the image piece \(\bar g(A)\) is well defined
independently of the choice of sticker in \(A\), and the equivariance
equation follows. The equation is preserved under composition and
inversion.

A full cap signature specifies every strict side choice in the cut
arrangement. The corresponding chamber is an intersection of open
half-spaces with \(\operatorname{int}\widetilde P\), so it has at most
one connected component. Thus two different chamber positions cannot
have the same complete signature. A hosting patch is an intersection of
the chamber closure with a facet, with its lower-dimensional boundaries
removed; the cut construction gives at most one such patch per hosting
facet.

Consequently \((M(A),h)\) addresses a sticker slot and \(\rho\) forgets
\(h\). This justification depends on the complete cut-chamber model. In
an arbitrary grip-set representation, signature injectivity would
require a separate hypothesis.

For a global geometric symmetry \(k\), cap and host labels are permuted:
\begin{equation}
M(kA)=kM(A),\qquad
\operatorname{Host}(kA)=k\operatorname{Host}(A).
\end{equation} A local twist need not act on every piece signature by
one common permutation of all cap labels; the destination signature is
determined by the actual destination chamber.

The two numerical formulas are the definitions of application rank and
sticker multiplicity. They count memberships and hosting patches.
Physical cut chambers have four-dimensional interiors, and sticker
interiors are three-dimensional, while the retained ranks range up to
\(19\). Thus rank cannot be either dimension or codimension in the
ambient four-space.

Nor does rank identify a move orbit: for example, O04 and O05 have the
same rank and sticker multiplicity but are distinct retained orbits.
Twenty corner stickers also do not imply a cyclic orientation variable
of order twenty; the specified corner frame group is \(A_5\).
\(\square\)

\subsection{Lemma 7: Transported frames and the ambient wreath
product}\label{lemma-7-transported-frames-and-the-ambient-wreath-product}

\subsubsection{Statement}\label{statement-6}

Let \begin{equation}
G=\langle H_c,T_c:c\in\mathcal C\rangle\le S_{\mathcal S}.
\end{equation} Under R, coherent transported frame choices give an
injective homomorphism \begin{equation}
G\hookrightarrow
\prod_{o=0}^{34}\left(H_o^{N_o}\rtimes S_{N_o}\right).
\end{equation} With chronological multiplication, the coordinates
satisfy \begin{equation}
(p,a)(r,b)=(r\circ p,c),
\qquad c_i=a_i b_{p(i)}.
\end{equation}

\subsubsection{Proof}\label{proof-6}

Fix one moving orbit. Choose a reference position and a reference
ordering of its stickers. Choose a legal transporter from the reference
position to every other position. These transporters identify all
sticker fibers with a common reference fiber.

For any legal map taking position \(i\) to position \(j\), composing
with the chosen transporters gives a return map on the reference fiber.
Its faithful permutation action belongs to the retained local frame
group \(H_o\). This constructs coherent coordinates; arbitrary
independent screen orderings would not automatically have this property.

Write the position permutation as \(p\), and the frame increment at
source position \(i\) as \(a_i\in H_o\). In the right-action convention,
an abstract frame state transforms by \begin{equation}
(i,h)\longmapsto(p(i),ha_i).
\end{equation} Applying \((p,a)\) and then \((r,b)\) gives
\begin{equation}
(i,h)\longmapsto
\bigl(r(p(i)),\,h a_i b_{p(i)}\bigr).
\end{equation} This proves the multiplication law, including its
reindexing.

Every legal move preserves each of its piece orbits, so the construction
applies independently to the \(35\) moving orbits. If all resulting
coordinates are trivial, every moving sticker is fixed, because each
fiber action is faithful. The center stickers are fixed by R. Hence the
original sticker permutation is the identity, proving injectivity.

Changing the coherent reference frames conjugates the coordinate
realization; it does not enlarge the legal group.

For a source-to-destination sticker permutation \(g\), a labelled state
array obeys \begin{equation}
L'[g(s)]=L[s].
\end{equation} This is consistent with the chronological convention used
throughout. \(\square\)

\subsection{Lemma 8: Necessary invariants and the conditional upper
bound}\label{lemma-8-necessary-invariants-and-the-conditional-upper-bound}

\subsubsection{Statement}\label{statement-7}

Under R and I, \begin{equation}
|G|
\le
\frac{\displaystyle\prod_{o=0}^{34}N_o!\,|H_o|^{N_o}}
{2^{43}5^2}.
\end{equation} The denominator is the exact index of the stated
invariant kernel in the ambient product. Equality with \(|G|\) is not
asserted.

\subsubsection{Proof}\label{proof-7}

Let \begin{equation}
H_o^{\mathrm{ab}}=H_o/[H_o,H_o],
\qquad
\chi_o:H_o\to H_o^{\mathrm{ab}}.
\end{equation} Write the abelian quotient additively and define
\begin{equation}
\Phi_o(p,a)=\sum_i\chi_o(a_i).
\end{equation} Lemma 7 gives \begin{equation}
\begin{aligned}
\Phi_o((p,a)(r,b))
&=\sum_i\chi_o(a_i b_{p(i)})\\
&=\sum_i\chi_o(a_i)+\sum_i\chi_o(b_{p(i)})\\
&=\Phi_o(p,a)+\Phi_o(r,b).
\end{aligned}
\end{equation} Thus \(\Phi_o\) is a homomorphism. Positional sign is
also a homomorphism because \begin{equation}
\operatorname{sgn}(r\circ p)=\operatorname{sgn}(r)\operatorname{sgn}(p).
\end{equation} Hypothesis I says that every generator lies in the
kernels of these maps, so every legal word does.

The local abelianizations are \begin{equation}
C_2^{\mathrm{ab}}=C_2,\quad
C_5^{\mathrm{ab}}=C_5,\quad
D_5^{\mathrm{ab}}=C_2,\quad
A_5^{\mathrm{ab}}=1.
\end{equation} For \(D_5\), the quotient by \(\langle R\rangle\) is
\(C_2\), while \begin{equation}
[R,F]=R^2
\end{equation} generates \(\langle R\rangle\). Thus
\([D_5,D_5]=\langle R\rangle\).

For completeness, \(A_5\) is perfect: take \begin{equation}
u=(1\,2)(4\,5),\qquad v=(2\,3)(4\,5).
\end{equation} Their commutator is a three-cycle. The derived subgroup
is normal, and all three-cycles are conjugate within \(A_5\). Since
three-cycles generate \(A_5\), its derived subgroup is all of \(A_5\).

There are thirty-five positional sign constraints, seven \(C_2\)
orientation constraints, one \(D_5\) orientation constraint, and two
\(C_5\) constraints.

These constraints are independent in the ambient product. Positional
parity can be prescribed using a transposition and identity frame
increments. Each orientation quotient value can be prescribed using one
frame coordinate and the identity positional permutation. Choices in
different orbits are independent. Therefore the joint invariant map is
onto a group of order \begin{equation}
2^{35}\cdot2^7\cdot2\cdot5^2=2^{43}5^2.
\end{equation} Its kernel has precisely this index. Lemma 7 and
hypothesis I place \(G\) inside that kernel, proving the bound.

A reflection-type element of the natural five-point \(D_5\) action
exchanges two pairs and is even as a sticker permutation. Its
abelianization bit therefore cannot be replaced by ordinary sticker
sign. Likewise, an odd positional permutation in one orbit cannot cancel
an odd positional permutation in another under I.

The bound leaves open additional restrictions and reachability
sufficiency. \(\square\)

\subsection{Lemma 9: Pure stars generate independent alternating
actions}\label{lemma-9-pure-stars-generate-independent-alternating-actions}

\subsubsection{Statement}\label{statement-8}

Under R and C, \(G\) contains a subgroup acting as \begin{equation}
A_{3600}\times A_{2400}\times A_{7200}^{\,7}
\end{equation} on the nine selected single-sticker orbits and fixing
every other sticker.

\subsubsection{Proof}\label{proof-8}

Words execute from left to right. Let \begin{equation}
t_C=(A\,B\,C)
\end{equation} be a relocated pure controller. If \(S\) fixes \(A,B\)
and sends \(C\) to \(X\), then the chronological word \begin{equation}
S^{-1}t_CS
\end{equation} has ordinary functional action
\(S\circ t_C\circ S^{-1}\), so it is exactly \begin{equation}
t_X=(A\,B\,X).
\end{equation} Conjugating the complete pure permutation gives a
complete pure permutation. Temporary collateral during its legal witness
is irrelevant to this net support statement.

For distinct nonbuffer positions \(X,Y\), direct evaluation gives
\begin{equation}
t_Y^{-1}t_X=(A\,Y\,X).
\end{equation} Together with the original stars and their inverses,
these supply every three-cycle containing \(A\). For distinct
\(x,y,z\ne A\), \begin{equation}
(A\,x\,y)(A\,z\,x)=(x\,y\,z).
\end{equation} Thus the stars generate every three-cycle.

Every even permutation is a product of three-cycles. For example,
express it as an even number of transpositions and pair them; a pair
sharing a point is a three-cycle, identical transpositions cancel, and a
disjoint pair can be rewritten as \begin{equation}
(a\,b)(c\,d)=(a\,b\,c)(a\,d\,c)
\end{equation} in the chronological convention. Conversely, every star
is even. The supported subgroup generated on this orbit is therefore
exactly \(A_N\).

The nine supported subgroups have disjoint full sticker supports. They
commute, and the multiplication map from their direct product is
injective: restricting an identity product to one selected orbit makes
that factor the identity. This proves the claim.

This identifies a supported subgroup of \(G\). It does not identify the
full legal group or its actions on the other twenty-six moving orbits.
\(\square\)

\subsection{Lemma 10: Balanced-color lower
bound}\label{lemma-10-balanced-color-lower-bound}

\subsubsection{Statement}\label{statement-9}

Under R and C, the set \(\Omega\) of reachable fixed-frame face-color
states satisfies \begin{equation}
|\Omega|\ge B,
\end{equation} where \begin{equation}
B=
\frac{3600!}{(6!)^{600}}\,
\frac{2400!}{(4!)^{600}}\,
\left[\frac{7200!}{(12!)^{600}}\right]^7.
\end{equation} These states can be attained with every unselected orbit
solved at macro boundaries.

\subsubsection{Proof}\label{proof-9}

Consider one selected orbit with \(N=600m\) positions and \(m\) copies
of each color. The stabilizer of its solved coloring in \(S_N\) is
\begin{equation}
Q=(S_m)^{600}.
\end{equation} Since \(m\in\{4,6,12\}\), this stabilizer contains a
transposition of two equal-colored positions and hence an odd
permutation. Therefore \begin{equation}
|A_N\cap Q|=\frac{|Q|}{2}.
\end{equation} The number of colorings in the \(A_N\) orbit is
\begin{equation}
\frac{|A_N|}{|A_N\cap Q|}
=
\frac{N!/2}{(m!)^{600}/2}
=
\frac{N!}{(m!)^{600}}.
\end{equation}

The odd same-color transposition is used only to describe the coloring
stabilizer. It need not itself be a legal move.

Lemma 9 gives independent alternating actions with full purity, so the
nine factors multiply and every unselected sticker remains solved at the
ends of the resulting macros. This gives \(B\).

A labelled state distinguishes every home sticker identity and is
determined by an element of \(G\). A face-color state forgets
distinctions between stickers with the same home-cell symbol. Thus \(B\)
is a lower bound for colored states, whereas Lemma 8 is an upper bound
for labelled states.

The colors here are \(600\) distinct mathematical symbols in a fixed
spatial frame. Repeated display RGB values or additional quotienting by
spatial and symbol symmetries would define a different state set and
require a separate count. \(\square\)

\subsection{Lemma 11: Word-ball and total-variation
bounds}\label{lemma-11-word-ball-and-total-variation-bounds}

\subsubsection{Statement}\label{statement-10}

Under R and C, use the unit-cost alphabet \begin{equation}
\mathcal A=\{H_c,T_c,T_c^{-1}:c\in\mathcal C\},
\qquad q=1800.
\end{equation} Some reachable face-color state has distance at least
\(46648\) from solved.

If \(\mu_{1000}\) is any distribution produced by \(1000\) successive
letters of this alphabet from solved and \(\nu\) is uniform on
\(\Omega\), then \begin{equation}
\|\mu_{1000}-\nu\|_{\mathrm{TV}}
>1-10^{-148595}.
\end{equation}

\subsubsection{Proof}\label{proof-10}

At most \(q^j\) words have length \(j\). Different words may give the
same state, so \begin{equation}
|\{x:d(\mathrm{solved},x)\le d\}|
\le
\sum_{j=0}^dq^j
=
\frac{q^{d+1}-1}{q-1}.
\end{equation}

The numerical comparisons needed here are the exact integer inequalities
\begin{equation}
1800^{46648}-1<1799B,
\label{eq:blueprint-1}
\end{equation} and \begin{equation}
1800^{1000}\,10^{148595}<B.
\label{eq:blueprint-2}
\end{equation} They can be evaluated without logarithms: form the
factorials in \(B\) as integer products, perform the indicated exact
multinomial divisions, and compare positive integers. These comparisons
are confirmed by the accompanying exact-integer audit. As an arithmetic
cross-check, \begin{equation}
1800^{46649}-1\ge1799B.
\end{equation} The cross-check says only when this counting argument
stops excluding coverage.

By Equation\textasciitilde{}\eqref{eq:blueprint-1}, the radius-\(46647\)
ball contains fewer than \(B\) states. Lemma 10 gives at least \(B\)
reachable states, so some state lies outside that ball.

Let \(A\) be the support of \(\mu_{1000}\). Then \begin{equation}
|A|\le1800^{1000},\qquad \mu_{1000}(A)=1.
\end{equation} Using the definition \begin{equation}
\|\mu-\nu\|_{\mathrm{TV}}
=\max_E|\mu(E)-\nu(E)|,
\end{equation} we obtain \begin{equation}
\begin{aligned}
\|\mu_{1000}-\nu\|_{\mathrm{TV}}
&\ge1-\frac{|A|}{|\Omega|}\\
&\ge1-\frac{1800^{1000}}{B}\\
&>1-10^{-148595},
\end{aligned}
\end{equation} where strictness follows from
Equation\textasciitilde{}\eqref{eq:blueprint-2}. No stationarity,
independence, or uniform choice of letters is needed for this support
argument.

These are statements about a specified alphabet and fixed-frame
face-color states. Macro applications, simultaneous native commands, and
camera motions are not single letters in this metric.

The distance bound is existential, not a lower bound for every scramble.
A known \(1000\)-letter scramble has an inverse solution of at most
\(1000\) letters because \(\mathcal A\) is inverse closed.

If a valid finite symmetry group \(L\) acts on the chosen state set,
quotienting may reduce counts: generally \begin{equation}
|\Omega/L|\ge\frac{|\Omega|}{|L|},
\end{equation} rather than preserving \(|\Omega|\). A quotient's goals
and metric must also be specified before transferring distance claims.
\(\square\)

\subsection{Lemma 12: Geometric normalization of the legal
group}\label{lemma-12-geometric-normalization-of-the-legal-group}

\subsubsection{Statement}\label{statement-11}

Under R, F, and E, the geometric action of \(K\) normalizes \(G\), its
intersection with \(G\) is trivial, and \begin{equation}
\langle G,K\rangle\cong G\rtimes K.
\end{equation}

\subsubsection{Proof}\label{proof-11}

Let \(k\in K\). Orthogonality and equivariance of poles give
\begin{equation}
k(U_c)=U_{kc}.
\end{equation} Let \(R_{c,a}\) denote a legal cap rotation corresponding
to \(a\in A_4\): it rotates the cap and fixes its complement. On the
complete subdivision, \begin{equation}
kR_{c,a}k^{-1}=R_{kc,kak^{-1}}.
\end{equation} The restriction \(kak^{-1}\) fixes the new cap axis and
is a proper rotation on its perpendicular three-space. This remains true
when \(\det k=-1\), because determinant is preserved by conjugation. It
is consequently an element of the new tetrahedron's \(A_4\) rotation
group.

By F and Lemma 6, every such local rotation is a finite word in
\(H_{kc},T_{kc}\). Therefore \begin{equation}
kGk^{-1}\subseteq G.
\end{equation} Applying the same argument to \(k^{-1}\) gives equality.

Now suppose \(g=k\) as sticker permutations, with \(g\in G\) and
\(k\in K\). Every center slot is fixed by \(g\), hence by \(k\).
Hypothesis E makes this mean \begin{equation}
kn_c=n_c\qquad\text{for every }c.
\end{equation} The poles span \(\mathbb R^4\). Otherwise a nonzero
vector orthogonal to every pole would give an unbounded line in the
intersection of the facet half-spaces, contradicting boundedness of
\(\widetilde P\). Thus \(k\) fixes a spanning set and is the identity.
Since \(G\) is a subgroup of the sticker symmetric group, its action is
faithful, and \(g\) is also the identity.

Normalization and trivial intersection imply that every element of
\(\langle G,K\rangle\) has a unique form \(gk\). Multiplication is
governed by the conjugation action of \(K\) on \(G\), giving the
asserted internal semidirect product.

The faithfulness requirements are essential. An action on a display
quotient that identifies centers, or a geometric representation without
a distinguished spanning center set, would not justify the intersection
argument. \(\square\)

\subsection{Distinct symmetry and motion
groups}\label{distinct-symmetry-and-motion-groups}

The groups appearing above serve different purposes.

{\def\LTcaptype{none} % do not increment counter
\begin{longtable}[]{@{}
  >{\raggedright\arraybackslash}p{(\linewidth - 2\tabcolsep) * \real{0.5000}}
  >{\raggedright\arraybackslash}p{(\linewidth - 2\tabcolsep) * \real{0.5000}}@{}}
\toprule\noalign{}
\begin{minipage}[b]{\linewidth}\raggedright
Group
\end{minipage} & \begin{minipage}[b]{\linewidth}\raggedright
Action and role
\end{minipage} \\
\midrule\noalign{}
\endhead
\bottomrule\noalign{}
\endlastfoot
\(SO(4)\) & Continuous orientation-preserving changes of
four-dimensional viewing frame \\
\(\operatorname{Sp}(1)\times\operatorname{Sp}(1)\) & Double cover of
\(SO(4)\), with diagonal \(\{\pm(1,1)\}\) kernel \\
\(K^+\) & The \(7200\) orientation-preserving symmetries preserving the
finite polytope \\
\(K\) & All \(14400\) geometric symmetries of the polytope \\
\(A_4\) at one cap & Its twelve local tetrahedral rotations \\
\(G\) & The finite group of legal permutations of the retained sticker
slots \\
\end{longtable}
}

To see the continuous covering directly, take \(A\in SO(4)\) and put
\(p=A(1)\). The map \(L_{p^{-1}}\circ A\) fixes \(1\) and restricts to
an element of \(SO(3)\) on the imaginary quaternions. Every such
three-dimensional rotation is \(x\mapsto qxq^{-1}\) for a unit
quaternion \(q\), so \(A\) has a unit-quaternion-pair representation.
The kernel calculation is the same as in Lemma 4.

A generic camera rotation does not preserve the finite set \(V\). A
local cap turn generally is not a global rigid symmetry. Under C, the
lower bound \(B>10^{151851}\) alone already separates the scale of the
legal action from the finite geometric group.

\subsection{Lemma 13: Conditions for sound symmetry
reduction}\label{lemma-13-conditions-for-sound-symmetry-reduction}

\subsubsection{Statement}\label{statement-12}

Normalization of \(G\) by \(K\) does not imply preservation of the
chosen \(1800\)-symbol word metric.

A symmetry quotient preserves shortest-path distance to a goal set if
the symmetries act by cost-preserving automorphisms of the actual
allowed state graph, preserve the goal set and the required protection
conditions, and quotient paths retain sufficient information to lift to
actual labelled paths.

\subsubsection{Proof}\label{proof-12}

Normalization says that a conjugated generator is an element of \(G\).
It says nothing about its length in \(\mathcal A\).

Already in one tetrahedral group, take \begin{equation}
H=(1\,2)(3\,4),\qquad T=(1\,2\,3).
\end{equation} Conjugating \(H\) by \(T\) gives a different half-turn.
This is none of the three chosen local symbols \(H,T,T^{-1}\), although
it is a finite word in them. Thus the local generating set is not
invariant under all local conjugations. Global geometric symmetry
supplies no automatic remedy for this failure. A separate audit would be
needed to show preservation of the particular global alphabet and its
primitive costs.

Now let \(\Gamma\) be the actual allowed state graph, with edge costs,
and let \(L\) act by cost-preserving graph automorphisms. Suppose the
goal set is \(L\)-invariant. The graph must already encode the relevant
allowed states, protected buffers or frames, legal generators, and any
restrictions on intermediate behavior.

Projection sends every path to a quotient path of the same cost.
Conversely, suppose an edge of a quotient path is represented by
\(u\to v\), while the current lifted state is \(x=\ell u\). The edge
lifts to \begin{equation}
x=\ell u\longrightarrow\ell v
\end{equation} with the same cost, because \(\ell\) is a graph
automorphism. Repeating this lifts the entire path. If its quotient
endpoint is a goal orbit, the lifted endpoint is a goal because the goal
set is invariant. Therefore the quotient and original minimum costs
agree.

To implement the lift, one must retain a representative edge, its actual
move, and the symmetry needed to align its source with the current
lifted state. Ordered frames, labels, and buffer identities must be
transported as part of this information.

For fixed-frame solved colors, a geometric permutation of positions
alone need not preserve the solved coloring or even the reachable color
set: it moves centers with distinct colors. A combined position action
and matching permutation of home-cell symbols may preserve solved, but
that is a separately defined action.

Likewise, setwise preservation of two buffers is insufficient when the
task protects them individually with ordered frames. The acting subgroup
must preserve the actual requirement or carry explicitly tracked buffer
data.

Without these conditions, a quotient may change the goal, admit
forbidden transitions, reduce costs artificially, or lose the data
needed to reconstruct a legal labelled solution. \(\square\)

\subsection{Lemma 14: Complete endpoint transport and certified
caching}\label{lemma-14-complete-endpoint-transport-and-certified-caching}

\subsubsection{Statement}\label{statement-13}

Let \(W\) and \(S\) be verified legal sticker permutations, with
chronological witnesses. Suppose the complete nonfixed map of \(W\) is
\begin{equation}
E_W=\{(u,W(u)):W(u)\ne u\}.
\end{equation} Then the complete nonfixed map of the chronological
conjugate \(S^{-1}WS\) is \begin{equation}
E_{S^{-1}WS}
=
\{(S(u),S(W(u))):(u,W(u))\in E_W\}.
\end{equation} Thus transporting both endpoints computes exactly the
same full permutation as replaying the conjugated legal witness,
including all collateral.

\subsubsection{Proof}\label{proof-13}

As an ordinary function, the chronological word \(S^{-1}WS\) is
\begin{equation}
S\circ W\circ S^{-1}.
\end{equation} For each \(u\), \begin{equation}
(S\circ W\circ S^{-1})(S(u))=S(W(u)).
\end{equation} Since \(S\) is bijective, \(S(W(u))\ne S(u)\) exactly
when \(W(u)\ne u\). Every moved source of the conjugate is therefore
represented, and no fixed source is included.

This proves equality on every sticker, not merely on a target orbit.
Collateral cycles are transported by exactly the same formula.

Consequently a cached full map can be certified by a complete seed map,
a verified setup witness, the endpoint transport calculation, and
consistent permutation conventions. Cache identity must distinguish the
underlying model, seed, and setup data sufficiently to avoid
substituting a different map.

The result licenses exact reuse of a permutation. It gives no numerical
performance guarantee and does not shorten the retained primitive
witness. \(\square\)

\subsection{Lemma 15: Guarded setup graphs and legal
inverses}\label{lemma-15-guarded-setup-graphs-and-legal-inverses}

\subsubsection{Statement}\label{statement-14}

Fix two buffer positions with their sticker frames pointwise protected.
On the finite graph of ordered target sticker frames, enlarging the
positive guarded alphabet by legal inverses cannot increase any minimum
unit-cost setup length.

For nonnegative unequal edge costs, minimum cost must be computed by a
suitable weighted shortest-path algorithm rather than ordinary
breadth-first search.

\subsubsection{Proof}\label{proof-14}

A vertex records a complete ordered target frame \begin{equation}
f=(s_1,\ldots,s_{k_o})
\end{equation} at a nonbuffer position. An allowed primitive \(g\)
induces the transition \begin{equation}
f\longmapsto(g(s_1),\ldots,g(s_{k_o})).
\end{equation} Allowed moves fix both buffer sticker fibers pointwise.
In the retained cap guard, forbidding every cap in \begin{equation}
M(A)\cup M(B)
\end{equation} is a sufficient condition for this fixed-buffer property.

Let \(E_+\) be the edges arising from allowed positive moves. Adding
their legal inverses gives an edge set \(E_{\pm}\) with \begin{equation}
E_+\subseteq E_{\pm}.
\end{equation} Every positive setup path is still available with the
same unit cost. Thus, with infinite distance allowed for unreachable
vertices, \begin{equation}
d_{\pm}(f_0,f)\le d_+(f_0,f).
\end{equation}

In this setting inverse moves also fix the buffers. Moreover,
\(H_c^{-1}=H_c\) and \(T_c^{-1}=T_c^2\), so inverse transitions already
have positive legal witnesses. Adding inverse letters can shorten paths
in the chosen alphabet without needing to enlarge its reachable
component.

For unit costs, breadth-first search explores vertices in nondecreasing
path length and finds minimum setup lengths in that graph. If costs are
nonnegative but unequal, path length is no longer path cost. A weighted
algorithm, such as Dijkstra's algorithm, is appropriate. Adding inverse
edges still cannot increase minimum cost provided old edges retain their
costs.

These statements compare precisely specified graphs. They do not claim
strict improvement, numerical speedup, or global optimality among
different buffer choices or setups allowed to move the buffers
temporarily. \(\square\)

\subsection{Scope tests: collateral, commutation, and quotient
equality}\label{scope-tests-collateral-commutation-and-quotient-equality}

Two elementary examples expose invalid stronger claims.

First, on six stickers the permutations \begin{equation}
u=(1\,2\,3),
\qquad
v=(1\,2\,3)(4\,5\,6)
\end{equation} have identical restrictions to the target orbit
\(\{1,2,3\}\) but different complete actions. Both are even. Equality in
a target-orbit quotient therefore does not justify replacing a
full-state macro. One must prove equality on the whole state, or
explicitly justify the collateral difference under the application's
current conditions.

Second, consider \begin{equation}
a=(1\,2)(5\,6),
\qquad
b=(3\,4)(6\,7).
\end{equation} If only stickers \(1,2,3,4\) are displayed, their visible
supports are disjoint. But their full actions do not commute: applying
\(a\) then \(b\) sends \(5\) to \(7\), whereas applying \(b\) then \(a\)
sends \(5\) to \(6\).

By contrast, permutations with disjoint full supports commute. Each
permutation fixes the support of the other pointwise, and a permutation
maps its own support to itself; checking a point in either support or
outside both gives equality of the two compositions.

These examples refute the stronger quotient and visible-support claims.
They are consistent with, and explain the necessity of, the
complete-support hypotheses in Lemmas 9 and 14.

\subsection{Theorem: Full 600-cell theory
package}\label{theorem-full-600-cell-theory-package}

\subsubsection{Statement}\label{statement-15}

Let \(V,P,K,K^+\), the pole convention, and hypotheses R, F, I, C and E
be as defined above. Then the following assertions hold, with the
hypotheses indicated separately.

\begin{enumerate}
\def\labelenumi{\arabic{enumi}.}
\item
  \textbf{Unconditional geometry and topology.} The radius-two polytope
  \(P\) is the regular \(\{3,3,5\}\) 600-cell with edge length
  \(2/\phi\), face vector \((120,720,1200,600)\), tetrahedral facets and
  icosahedral vertex links. Five facets meet at each edge and twenty at
  each vertex. Its boundary is radially homeomorphic to \(S^3\), with
  integral homology \(\mathbb Z\) in degrees zero and three, zero in
  degrees one and two, and trivial fundamental group. Its polar is a
  regular 120-cell. The finite checker verifies its specified arithmetic
  and incidence statements; mod-2 Betti numbers alone do not establish
  this topological identification.
\item
  \textbf{Unconditional geometric symmetry.} The unit vertices \(V/2\)
  form \(I^*\), and the left/right quaternion action identifies \(K^+\)
  with \((I^*\times I^*)/\langle(-1,-1)\rangle\), of order 7,200.
  Quaternion conjugation extends this action to \(K=W(H_4)\) of order
  14,400, with the Coxeter presentation stated in Lemma 4. These finite
  groups are distinct from \(SO(4)\), the rotation group of a cap and
  the legal twist group.
\item
  \textbf{Retained cut model, under R and F.} The 600 oriented poles
  describe 300 antipodal axes with two outer layers at normalized
  distance \(121/125\) and a fixed middle layer during an outer-layer
  turn about that axis. The retained model has 259,800 sticker slots,
  177,120 surface pieces, 600 fixed centers and the stated 35 moving
  orbits. Sticker-to-piece incidence is equivariant under legal twists.
  Complete cap signatures and hosting sets provide the specified
  addresses, while \(|M|-1\) is a membership rank rather than a
  dimension. The pair \(H_c,T_c\) generates the full local \(A_4\)
  action. The count of 433 sticker regions per facet remains a
  retained-model input.
\item
  \textbf{Legal actions and conditional bounds.} Under R and F, the
  legal group embeds in the stated product of transported-frame wreath
  products, with chronological multiplication
  \((p,a)(r,b)=(r\circ p,c)\) and \(c_i=a_i b_{p(i)}\). With the
  additional generator-invariant hypothesis I, positional parity and
  orientation abelianization give

  \begin{equation}
  |G|\leq \frac{\prod_o N_o!\,|H_o|^{N_o}}{2^{43}5^2}.
  \end{equation}

  With the separate pure-controller hypothesis C, the supported action
  contains \(A_{3600}\times A_{2400}\times A_{7200}^{7}\) and yields the
  balanced fixed-frame color lower bound \(B\) of Lemma 10. Neither
  bound determines the exact order of \(G\) or proves sufficiency of the
  listed invariants.
\item
  \textbf{Specified word metric, under R and C.} For the 1,800-symbol
  alphabet \(\{H_c,T_c,T_c^{-1}\}\), the word-ball inequality implies a
  worst-case distance of at least 46,648. Any 1,000-step distribution
  from solved has total-variation distance strictly greater than
  \(1-10^{-148595}\) from uniform reachable fixed-frame face-color
  states. This is not a lower bound for each scramble; a recorded
  1,000-letter word has an inverse of at most 1,000 letters. Labelled
  states, face-color states and symmetry quotients remain distinct.
\item
  \textbf{Extended geometric action, under R, F and E.} The faithful
  geometric action on the complete retained subdivision normalizes
  \(G\). Since legal moves fix every distinguished center and those
  centers span the geometric space, \(G\cap K=\{1\}\) and
  \(\langle G,K\rangle\cong G\rtimes K\). Normalization does not
  automatically preserve the chosen move alphabet or cost. Sound
  shortest-path reduction requires the graph, goal, protection and
  reconstruction conditions of Lemma 13.
\item
  \textbf{Sound optimization.} Under the complete-action and
  guarded-graph hypotheses of Lemmas 14 and 15, transport of both
  support endpoints computes the full conjugated net permutation, and
  adding legal inverse edges cannot increase minimum unit-cost setup
  length. Unequal nonnegative costs require weighted shortest-path
  search. Equality on a target orbit does not establish equality of
  full-state macros; disjoint full supports suffice for commutation,
  while disjoint visible supports do not. No numerical speedup, new
  solve or performance record follows from these statements.
\end{enumerate}

\subsubsection{Proof}\label{proof-15}

Lemmas 1--3 establish the unconditional coordinate geometry, face
incidences, radial homeomorphism, integral homology, fundamental group,
and polar duality required in part (1). The finite-checker discussion
specifies the exact enumeration algorithm, the supporting-facet
completeness argument, and the conclusions outside its computational
scope.

Lemma 4 identifies the quaternion symmetry action, proves its kernel,
establishes both the lower and upper symmetry counts, and proves the
Coxeter presentation required in part (2). The continuous covering
argument and the group-role table distinguish camera motions, finite
geometric symmetries, local cap rotations, and legal sticker
permutations.

Lemma 5 gives the exact scale conversion for part (3). The
retained-model definitions specify the \(600\) oriented caps, their
\(300\) antipodal axes, signatures, hosting sets, and census without
claiming an independent cut reconstruction. Lemma 6 proves local \(A_4\)
generation and piece equivariance and explains the application rank.

Lemma 7 proves the faithful ambient embedding and chronological frame
law. Lemma 8 derives the conditional invariant upper bound under I,
including independence of its ambient index factors. Lemmas 9 and 10
derive the independent alternating actions and balanced-color lower
bound under the separate pure-controller hypothesis C. These establish
part (4) without asserting full reachability or an exact legal-group
order.

Lemma 11 proves part (5) from the word count, the color lower bound, and
explicit integer inequalities. It distinguishes worst-case existence
from individual scramble length and distinguishes fixed-frame color
states from labelled and symmetry-quotiented states.

Lemma 12 proves part (6)'s normalization, trivial intersection, and
semidirect product under the complete geometric realization E. Its proof
uses the faithful sticker action and the spanning geometric center set
explicitly. Lemma 13 shows why normalization alone is insufficient for
the chosen metric and proves the conditions and reconstruction
obligation for sound shortest-path symmetry reduction.

Finally, Lemmas 14 and 15 prove the two optimization statements in part
(7). The complete-support examples establish the stated limitations on
macro substitution and commutation. None of these arguments asserts a
numerical speedup, a new solve or performance record, or a determination
of the unknown full legal-group order. \(\square\) \newpage

\section{From theory to
implementation}\label{from-theory-to-implementation}

The preceding results offer several ways to reduce work without changing
the puzzle. The practical target is preservation of the full labelled
action. A faster operation is useful only if it carries the same piece
identities, orientations, collateral effects and history semantics as
its reference implementation.

\subsection{Normalization without changing model
identity}\label{normalization-without-changing-model-identity}

For the radius-two coordinates, the supporting-facet calculation gives
squared inradius

\begin{equation}
r^2=\frac{(4+6\phi)^2}{16+24\phi}
=1+\frac32\phi=\frac{\phi^4}{2}.
\end{equation}

The retained pole convention has \(\|n_0\|^2=\phi^6+3=8\phi^2\). Thus
the ratio between its inradius and the radius-two model's inradius is
\(4/\phi\). A suitable orthogonal identification and that scale relate
the two geometric descriptions. They do not determine the application's
cell numbering or sticker IDs. The normalized cap plane at \(121/125\)
of the cell distance lies inward by \(4/125\) of the inradius. This is a
geometric distance ratio, not a statement that 3.2 percent of the pieces
move.

Keeping the normalization explicit prevents a particularly subtle error:
applying the numeric constant 0.968 to coordinates whose supporting
planes have a different distance from the origin. Any replacement
geometry pipeline should transform both positions and plane equations
coherently, then verify its cell and sticker correspondence against the
immutable model manifest.

\subsection{State representation and incremental
updates}\label{state-representation-and-incremental-updates}

Let \(S\) be the fixed set of 259,800 sticker slots. Store \(L[s]\) as
the identity of the sticker currently occupying slot \(s\). A legal
permutation \(m:S\to S\) is applied by

\begin{equation}
L_{\mathrm{new}}[m(s)]=L_{\mathrm{old}}[s].
\end{equation}

Let \(D=\operatorname{supp}(m)\), and let \(\pi:S\to\mathcal P\) map a
slot to its host piece position. Since a permutation preserves its own
support, only the destination positions in \(\pi(D)\) can have changed
their occupants or ordered frames. This licenses an incremental progress
update: update the inverse location map and recalculate correctness for
the affected piece positions, then adjust the per-orbit totals. It does
not license dropping stickers from the mechanical state.

The present \texttt{PuzzleState.refresh} reconstructs location and
correctness information with full-array work. A future incremental
implementation should retain the full refresh as an independent oracle.
For each legal generator, its inverse, a composite macro and an imported
replay, compare the complete inverse-location arrays, every
piece-correctness bit and every per-orbit count. A discrepancy must
reject the optimization before it can affect a saved session. Dense
moves may make a full vectorized refresh faster than bookkeeping over a
large support; the decision should be measured against support size.

The model identity and the full-state hash have different purposes. The
former fixes geometry, IDs, cuts, frames and generator meaning; the
latter identifies a particular labelled configuration. Neither may be
replaced by a smaller displayed subset. In particular, an incremental
noncryptographic XOR digest is not a substitute for the current SHA-256
boundary used by previews and certificates. A change to hashing
semantics would require an explicit protocol version and new validation.

\subsection{Certified net maps and
caches}\label{certified-net-maps-and-caches}

A long legal word may have a small net support. The endpoint-transport
result in the blueprint justifies caching the complete conjugated net
permutation while retaining its primitive witness. Applying that cache
can reduce execution work from repeated full-word replay to movement of
the net support. This is a distinction between execution cost and
mathematical word length: it does not shorten the recorded legal
witness.

A cache record should bind at least the model identifier, frame
convention, seed identifier and digest, setup word and direction,
complete net-map digest, and protected-orbit contract. Retain both
source and destination endpoints and verify bijectivity. An orbit-only
map is inadequate for a raw seed whose other effects occur in later
stages. Two macros with the same target three-cycle may be incompatible
substitutes because they move different collateral pieces.

There are also two separate symmetry-caching opportunities. For
geometric data, an exactly verified element of \(K\) can transport
incidence and sticker endpoints. For search, the same element is useful
only if it preserves the allowed transitions, edge costs, distinguished
buffers and goal condition, or if the search explicitly transforms all
of them. The theorem \(K\leq N(G)\) concerns legal reachability; it does
not establish that \(K\) acts by isometries of the chosen 1,800-symbol
metric.

\subsection{Search graphs, orientation and
cost}\label{search-graphs-orientation-and-cost}

The guarded setup graph tracks an \textbf{ordered target sticker frame},
with both buffers fixed. Forgetting frame order merges distinct
orientations. That merge can yield a setup that places the correct piece
in the correct position while leaving a nonabelian orientation error.
The retained controller graph has at most

\begin{equation}
V_o=(N_o-2)|H_o|
\end{equation}

candidate target frames for orbit \(o\). A transition that applies a
guarded generator must transport the whole frame. With \(b\) generators
and \(k\) sticker entries per frame, straightforward graph construction
costs \(O(V_obk)\), apart from generator storage. A predecessor tree
needs \(O(V_o)\) records rather than storing a complete setup word at
every vertex.

Adding legal inverses to a positive unit-cost alphabet cannot increase
any shortest setup distance. It may nevertheless increase the time spent
building the graph, so shorter words and lower build latency should be
measured separately. Ordinary BFS is appropriate for unit costs.
Primitive length, elapsed execution cost and ergonomic cost are
different nonnegative weight functions; a weighted search must use an
appropriate algorithm such as Dijkstra's algorithm. Bidirectional search
needs correctly reversed transitions and compatible fixed-buffer
conditions.

The retained workload tables identify O02, O15 and O29 as expensive
targets: together they account for about 40.34 percent of the
constructive primitive bound while requiring about 3.8 percent of its
stars. Shorter complete witnesses for those orbits could therefore
improve the primitive bound disproportionately. They need not improve
interactive speed by the same amount when the workbench already executes
cached net permutations. Every replacement requires full-support replay
and revalidation of the collateral order.

\subsection{Stable rendering through finite
symmetry}\label{stable-rendering-through-finite-symmetry}

The geometric relationship between cells suggests an instancing design.
Let \(v_j\) be vertices of one reference cell's subdivided geometry and
let \(R_c\) place that geometry in cell \(c\). Let \(Q\in SO(4)\) be a
camera rotation. A candidate rendering pipeline computes

\begin{equation}
v_{c,j}=R_cv_j,\qquad u_{c,j}=Qv_{c,j},\qquad
p_{c,j}=\Pi_3\bigl(\Pi_4(u_{c,j})\bigr).
\end{equation}

Here \(\Pi_4\) is a specified 4D-to-3D projection and \(\Pi_3\) the
final screen projection. For example, a perspective projection may use
\(\Pi_4(x,y,z,w)=d(x,y,z)/(d-w)\), restricted away from its singular
plane. Clipping must handle the valid domain before division; a small
denominator is not a reason to corrupt coordinates or alter the puzzle
state. The real application may use additional shrink, visibility and
view parameters, which must be reproduced as well.

The web mesh asset has 30,480 base vertices and 10,160 base triangles.
Replicating that mesh over 600 cells represents 6,096,000 triangles
before view-dependent rejection. This is a count for the web asset,
\textbf{not a measurement of the native DirectX renderer's mesh}.
Sticker count alone is therefore an inadequate graphics workload
descriptor. Profiling should report actual submitted vertices,
triangles, lines, draw calls and uploaded bytes for the native path.

One prospective GPU implementation would retain immutable reference
geometry and cell transforms, update the camera in a small uniform
buffer, and obtain colors and selection styles through a sticker-ID
lookup. Such a design could reduce repeated CPU projection and packing.
It still has to reproduce clipping, ordering, transparency and
selection. The installed Managed DirectX path and device capabilities
constrain what is practical; a renderer rewrite is a separate, tested
engineering change, not a consequence already implemented by the group
theory.

The existing native optimization reuses a 28-byte-per-vertex upload
array around a pinned upstream renderer. This reduces allocation and
marshaling while preserving the established projection path. It leaves
CPU projection and triangle/line packing in place. A useful next
experiment must therefore time those stages separately before
attributing the remaining delay to the GPU.

\subsection{Filtering, framework visibility and
picking}\label{filtering-framework-visibility-and-picking}

A filter is a predicate on pieces, potentially using their solved
identity, current position, orbit, rank, correctness or an explicit
context rule. It determines which geometry is displayed; it does not
define a smaller puzzle. All stickers belonging to a hidden piece must
be excluded consistently from rendering and hit testing. Buffer or
preview context outside the selected predicate should be explicit, so
additional pieces cannot be confused with a failed filter.

The 600-cell framework is a separate draw layer. Hiding it must suppress
its lines at every draw entry point, including animation and direct
redraws. It must retain the cell planes and clipping information used by
visible stickers. Removing structural clipping data to remove a visual
outline conflates geometry with presentation and can produce both
rendering and picking defects.

Render and pick operations should consume a coherent snapshot:

\begin{equation}
\mathcal R=(\text{revision},L,Q,F,\text{frame policy},\text{motion policy}).
\end{equation}

The picker's candidate set must use the same visibility predicate \(F\)
as the displayed scene. A cached projection is valid only for the
matching revision, camera and view parameters. After a sampled motion
frame, exact picking needs either a refreshed projection of the complete
visible geometry or a mathematically equivalent ray test against that
geometry. Invisible triangles and stale axes must not remain
interactive.

During movement, adaptive detail can improve responsiveness while the
complete state remains unchanged. When movement stops, restoring all
visible geometry is a distinct latency that should be reported. A
sampled frame must never be counted as a successful full-detail frame in
the 30 FPS qualification workload.

\subsection{Transactions and optimization
boundaries}\label{transactions-and-optimization-boundaries}

The state transition remains authoritative. A preview binds the current
revision, full-state hash and protection constraints; commit rechecks
all three. The complete net permutation is applied and journalled
atomically before the new head is exposed. Rendering consumes the
resulting immutable snapshot. Changing frame visibility, camera pose or
sampling policy must leave the state hash unchanged.

Reset and import require complete validation before durable writes and a
recoverable previous state. Checkpoints and log exports should record
the model identity, chronological moves or certified macros, and enough
state information to verify replay. A failed image update after a
successful database commit must not cause the operation to be applied
twice. These contracts matter especially for scrambling: color transport
follows the committed labelled permutation rather than a partly updated
display array.

The rendering queue may coalesce obsolete camera snapshots when a newer
view is available. It may not coalesce or discard committed mechanical
moves from the history. Likewise, progress reports may cache aggregates,
but sessions reports must describe durable operations rather than
estimates derived from animation completion.

\subsection{What performance evidence
supports}\label{what-performance-evidence-supports}

The frozen 0.2.4 validation recorded an optimized full-detail mean of
453.45 ms and a 95th percentile of 533.82 ms at a \(931\times604\)
viewport. The comparison path averaged 626.57 ms. The 27.63 percent
reduction comes from 18 fixed-pose redraw measurements, not a sustained
rotation test. The reciprocal of 453.45 ms is approximately 2.2 frames
per second; it is not a qualified 1080p frame rate.

For a mean baseline \(T\) and a measured nonaccelerated fraction \(f\),
an idealized Amdahl model gives

\begin{equation}
T_{\mathrm{new}}=fT+\frac{(1-f)T}{r}.
\end{equation}

Even infinite acceleration of the second term could reduce 453.45 ms to
33.33 ms only if \(f<0.0736\). The prior total process-CPU figure is not
a measurement of this critical-path fraction; CPU and GPU overlap,
waiting, packing and submission must be separated before using the model
quantitatively. More GPU throughput alone cannot be promised to satisfy
the target.

The target minimum remains \textbf{30 FPS at 1920 by 1080 with all
259,800 stickers present at full detail}. No specific minimum GPU model
is certified. Intel HD Graphics 620 supports the measured
filtered/adaptive workflow but has not passed this full-detail
qualification. Neither the new theory nor the exact geometry audit
changes that status.

{\def\LTcaptype{none} % do not increment counter
\begin{longtable}[]{@{}
  >{\raggedright\arraybackslash}p{(\linewidth - 2\tabcolsep) * \real{0.5000}}
  >{\raggedright\arraybackslash}p{(\linewidth - 2\tabcolsep) * \real{0.5000}}@{}}
\toprule\noalign{}
\begin{minipage}[b]{\linewidth}\raggedright
Proposed experiment
\end{minipage} & \begin{minipage}[b]{\linewidth}\raggedright
Evidence required before acceptance
\end{minipage} \\
\midrule\noalign{}
\endhead
\bottomrule\noalign{}
\endlastfoot
Incremental state/progress updates & Complete state, inverse-location
map and orbit counts agree with full recomputation \\
Symmetry or macro cache & Complete labelled net map equals legal
witness, including all collateral \\
Inverse or weighted setup search & Correct ordered frame, both buffers
fixed, measured cost and reconstructed legal word \\
GPU projection or instancing & Clipping and reference-image agreement,
stable sticker IDs and matching picking \\
Full-detail rotation qualification & At least 60 seconds each of warmed
3D and 4D motion at 1080p; full workload; frame-time distribution and
presentation method \\
Adaptive motion & Interaction latency, exact restoration latency,
unchanged mechanical hash and correct post-motion picking \\
\end{longtable}
}

For the full-detail qualification, report mean, median, 95th and 99th
percentile frame times; target a 95th percentile no greater than 33.33
ms as well as sustained 30 FPS. Identify the application build, driver,
viewport, clipping settings and actual visibility count. Check
representative images against the rendering reference, and verify
identical labelled hashes before and after camera-only runs. Cold
construction, restore latency and device recovery are separate
measurements. These are acceptance criteria for future work, not claims
that the present supplement ran those graphics experiments.

\section{Reproduction and
interpretation}\label{reproduction-and-interpretation}

The accompanying standard-library geometry audit is self-contained. It
constructs the 120 exact vertices, enumerates adjacency and simplex
cliques, verifies supporting facets, forms simplicial boundary maps over
\(\mathbb F_2\), computes graph distances, and enumerates the
quaternion-induced symmetry maps. Its JSON result records both numerical
outputs and the scope boundary. Running it does not open the application
or read a user session.

The original 35-orbit algorithms paper and the curated replay archive
remain unchanged. They supply the complete algorithm cards and finite
words. The existing configuration-bound and extended-derivation audits
remain the references for retained controller and arithmetic evidence.
The new checker addresses base geometry; it does not replace those
audits.

Rethlas was used at upstream revision
\texttt{887cc46427636bbdd235160a112f9a30ae81d040}, with separate
generation and verification invocations of Astra. The stock file-based
runner could not read its task inputs under the nested Windows execution
policy. The adaptation supplies the reviewed public inputs as text to
read-only agents and lets the external coordinator persist outputs
through Rethlas's memory and strict verification schema. The coordinator
retains rejected candidates and repair reports; acceptance requires the
verifier's error and gap lists both to be empty. Public records contain
the reviewed blueprint, exact verdicts, scope and hashes, while local
paths, raw execution logs and unrelated material are excluded.

The first complete inline blueprint, comprising fifteen lemmas and the
combined theorem, received the strict verdict \textbf{correct}, with
zero critical errors and zero gaps. No mathematical repair round was
required. Earlier file-policy attempts produced no candidate and are not
counted as completed reviews. The exact returned blueprint and verdict
are retained separately. This typeset exposition removes the
generation-time status paragraph, normalizes equation references, and
replaces imperative task wording in the combined statement with an
equivalent declarative statement. The engineering discussion is a
separately authored analysis.

The resulting mathematical review is informal. It neither proves the
exact legal-group order nor formally verifies the entire software stack.
Source provenance, exact finite checks and explicit hypotheses make its
conclusions inspectable; independent reproduction and mathematical
review remain possible without the private design transcript.

\newpage

\section{Reading guide}\label{reading-guide}

\begin{enumerate}
\def\labelenumi{\arabic{enumi}.}
\tightlist
\item
  \textbf{Hypercubing community, Formal Introduction to Grip Theory.}
  Defines active grips, attitudes and the distinction between general
  piece states and twist-reachable states. Read it alongside the
  blueprint's frame convention; the program's finite sticker action need
  not retain an invisible geometric attitude.
  \url{https://hypercubing.xyz/theory/grip-theory/formal/}
\item
  \textbf{Hypercubing community, Piece Invariants.} Background on parity
  and orientation constraints. Its hypercube monoflip example is
  illustrative rather than a full-600-cell certificate.
  \url{https://hypercubing.xyz/theory/invariants/}
\item
  \textbf{Hypercubing community, God's Number.} Explains metric-specific
  counting arguments and the distinction between a worst-case bound and
  a solution length for one state.
  \url{https://hypercubing.xyz/theory/gods-number/}
\item
  \textbf{Hypercubing community, Cut Depths.} Establishes normalized
  cut-distance conventions and illustrates how changing cuts changes a
  puzzle's pieces. It does not enumerate the retained full-600-cell
  subdivision. \url{https://hypercubing.xyz/theory/cut-depths/}
\item
  \textbf{John C. Baez, From the Icosahedron to E8.} Primary background
  for binary icosahedral quaternions and 600-cell coordinates;
  arXiv:1712.06436v2. \url{https://arxiv.org/abs/1712.06436}
\item
  \textbf{Jihyun Choi and Jae-Hyouk Lee, Binary Icosahedral Group and
  600-Cell.} \emph{Symmetry} 10(8), 326 (2018). A primary geometric
  treatment connecting quaternion multiplication to the regular
  polytope. \url{https://doi.org/10.3390/sym10080326}
\item
  \textbf{David Vogan, Regular polyhedra and Coxeter groups.} Tufts
  lecture, 25 January 2019. Background for flag counting and finite
  reflection groups. \url{https://math.mit.edu/~dav/regpolyTUFTSHO2.pdf}
\item
  \textbf{Rethlas.} Upstream generation and independent
  proof-verification workflow. The companion provenance describes the
  local I/O adaptation and records the actual review outcome.
  \url{https://github.com/frenzymath/Rethlas}
\item
  \textbf{C600 Studio research materials.} The technical report,
  unchanged original algorithms paper, synthetic replay supplement,
  audits, and Full Detail Rotation requirements.
  \url{https://github.com/KonomiYuzu01/C600-MPUlt/tree/main/research}
\item
  \textbf{Andrey Astrelin, Magic Puzzle Ultimate.} Original simulator
  and renderer underlying this mod.
  \url{https://superliminal.com/andrey/mpu/}
\end{enumerate}

Web references were checked on 7 September 2026. These links identify
public technical sources, not personal conversation histories.

\end{document}
